\chapter*{问题集}
以下历史疑点均已在相应年份的答案中核对并处理，保留本页作为复核索引。
\begin{enumerate}
	\renewcommand{\labelenumi}{\arabic{enumi}.}
	\item
	2010 年第一题（2）：算符平方根 $X=A^{1/2}$ 的定义是 $X^2=A$。一般并没有
	$A^{1/2}B^{1/2}=(AB)^{1/2}$；只有在 $A,B$ 对易且所选平方根分支相容等附加条件下才可这样合并。
	\item
	2011 年第三题：已重新计算简并微扰矩阵、能级分裂和相应本征态。
	\item
	2012 年第五题（3）：已用二重对易子完整计算变分参数与能量上限。
	\item
	2014 年第一题：过渡矩阵依赖两组基矢的排列和相位约定；固定约定后各答案等价。对于一般对易算符，应先对其中一个算符对角化，再在其每个简并子空间内对角化另一个算符。
	\item
	2014 年第二题：已用简并微扰论处理 $|12\rangle,|21\rangle$ 子空间，并写出非简并态及简并态的一级波函数修正通式。
	\item
	2014 年第四题（2）：已按平移后的谐振子精确时间演化重新计算旧基态的存留概率。
	\item
	2016 年第五题（3）：轨道态和自旋态的交换因子分别为 $(-1)^L$ 与 $(-1)^{S+1}$；电子总态反对称要求 $L+S$ 为偶数。
\end{enumerate}
